\documentclass[../main.tex]{subfiles}
\begin{document}
%==========================================TIÊU ĐỀ TẬP========================================
\begin{table}[h]
    \begin{adjustwidth}{-1cm}{}
        \begin{tabular}{>{\centering\arraybackslash}p{6cm}|>{\raggedright\arraybackslash}p{12.5cm}}
            \multirow{5}{6cm}
            {\\[-40pt]
                \flaregothic\fontsize{20pt}{20pt}\selectfont \centering
                \color{eptype}HISTOIRE\color{black}\\
                \freshbot\fontsize{72pt}{72pt}\selectfont
                \color{epnum}32\\[6pt]
                \cafeteria\fontsize{20pt}{20pt}\selectfont\color{epnum}(D-32)\color{black}
                \\[18pt]
                \color{black}
            }
             &
            {\fotskip\fontsize{18pt}{18pt}\selectfont \ruby{二}{ふた}つの\ruby{翼}{つばさ}と\ruby{一}{ひと}つの\ruby{道}{みち}}
            \\[6pt]
             & {\cafeteria\fontsize{18pt}{18pt}\selectfont Two Pairs of Wings, One Path}                                \\[4pt]
             & {\vnmsans\fontsize{18pt}{18pt}\selectfont Hai đôi cánh và một con đường}                                 \\[8pt]
             & {\caviar{\fontsize{13pt}{13pt}\selectfont Nhân vật: \emp{kanata} \emp{kson} \emp{kuon2} \emp{ryuuhou2}}}
            \\[8pt]
        \end{tabular}
    \end{adjustwidth}
\end{table}
%=========================================TRUYỆN CHÍNH========================================
\color{black}

\txtbx[2]{
    {\color{vol} \uline{Tóm tắt:}} Không còn là những chuyến tụ tập ồn ào của cả thế hệ, Kanata một mình bắt chuyến tàu hướng về ngoại ô Kawachi vào một buổi chiều cuối tháng Bảy, mang theo cuốn album ảnh đã ngủ yên suốt một năm trời. Trong căn phòng tầng bảy lộng gió, hai người bạn thân từng sống chung dưới một mái nhà cùng lật giở lại những ký ức từ thuở chập chững vào nghề: những buổi tập nhảy ướt đẫm mồ hôi, bữa ăn vụng trộm đêm muộn và cả những giọt nước mắt chia ly tưởng chừng không bao giờ hàn gắn. Cùng rảo bước qua những con dốc thanh bình của phường Anwa, họ giãi bày về những ngã rẽ tương lai—nơi một người tiếp tục bay cao trên bầu trời idol và một người tự do sải cánh kiến tạo không gian cho riêng mình. Dẫu hai đôi cánh chọn hai chân trời khác biệt, lời hứa năm nào vẫn vẹn nguyên, gắn kết họ trên cùng một hành trình tìm kiếm hạnh phúc.
}

\pov{Amane Kanata}

Chuyến tàu cao tốc tuyến Chuo hôm nay vắng lặng hơn thường lệ.

Tôi ngồi bên khung cửa kính, hai tay ôm chiếc túi xách bằng vải bố màu xanh lam đặt ngay ngắn trên đùi. Bên trong chiếc túi là một vật thể hình chữ nhật khá nặng, mép bìa bọc nhung xanh đã hơi sờn qua năm tháng.

Đó là cuốn album ảnh kỷ niệm của Thế hệ thứ Tư.

Suốt hơn ba trăm ngày qua, cuốn album ấy nằm lặng lẽ ở ngăn kéo dưới cùng của bàn làm việc trong căn hộ của tôi. Tôi đã không dám mở nó ra một lần nào. Mỗi khi nhìn thấy gáy sách quen thuộc, lồng ngực tôi lại thắt lại bởi một nỗi sợ vô hình: sợ rằng những nụ cười rạng rỡ in trên giấy bóng kia sẽ chỉ càng khắc sâu thêm sự trống trải của hiện tại.

Thế nhưng hôm nay, mọi chuyện đã khác.

Sau buổi chiều tái ngộ ở kho hàng ngoại ô và bữa tiệc nướng rộn rã trên sân thượng tầng bảy tuần trước, lớp sương mù dày đặc trong lòng tôi đã hoàn toàn tan biến. Tôi nhận ra mình không còn phải lén lút trốn chạy quá khứ nữa.

Và tôi muốn cùng cậu ấy lật lại từng trang ký ức này, bằng một nụ cười trọn vẹn nhất.

\ct

Bước xuống ga Kyuukami, cái nắng cuối tháng Bảy dát một màu vàng mật ong óng ả lên những mái nhà ống và những rặng cây phong xanh rì ven đường. Tôi men theo con dốc quen thuộc dẫn về ký túc xá nhà Hikawa.

Cánh cổng gỗ lớn khép hờ. Tiền sảnh tầng một mát mẻ và yên ắng, thoang thoảng mùi trà lúa mạch quen thuộc.

Anh Kuon đang ngồi bên bàn làm việc kiểm tra lại sổ sách và tài liệu quản lý. Nghe thấy tiếng bước chân rụt rè ở cửa, anh ngẩng đầu lên, nở nụ cười điềm đạm quen thuộc: ``A, Kana-chan đấy à? Hôm nay đi có một mình thôi sao?''

``Vâng ạ,'' tôi cúi đầu chào anh, khẽ ôm chặt chiếc túi vải. ``Hôm nay em được nghỉ lịch tập, muốn sang gặp riêng Kson một lát. Cậu ấy có ở nhà không anh?''

``Có chứ,'' anh với lấy bình trà lúa mạch rót cho tôi một ly đầy đá mát lạnh. ``Hôm nay bố vớiKaiho sang xưởng cơ khí bên kia thành phố làm, còn Kson sửa xong kênh phụ từ trưa thì đang ngồi nghe nhạc ở trên phòng tầng bảy đấy. Lên đi, con Rồng ấy vừa mới than thở là nhớ mấy đứa xong.''

Tôi nhận lấy ly trà, uống một ngụm mát rượi rồi mỉm cười: ``Cảm ơn Kuon-senpai ạ! Em lên với cậu ấy đây.''

``Ừ, cứ tự nhiên như ở nhà nhé,'' anh Kuon vẫy tay, ánh mắt đong đầy sự ấm áp của người anh cả luôn lặng lẽ che chở cho tất cả chúng tôi.

\ct

Tôi leo từng bậc cầu thang gỗ lên tầng bảy. Hành lang tầng trên cùng yên ắng đến mức chỉ nghe thấy tiếng ve sầu râm ran vọng lại từ tán cây phong ngoài cửa sổ.

Cánh cửa phòng hé mở một khe nhỏ. Từ bên trong, giai điệu mộc mạc của một bản acoustic vang lên nhè nhẹ.

Tôi gõ nhẹ ba tiếng lên cánh cửa gỗ.

``Cửa không khóa đâu, vào đi!'' giọng nói khàn khàn, sảng khoái quen thuộc cất lên.

Tôi đẩy cửa bước vào. Kson đang ngồi xếp bằng trên tấm thảm cói giữa phòng, trên người là chiếc áo phông đen rộng thùng thình in hình đầu lâu cùng chiếc quần ngố thoải mái. Mái tóc tím than dài được búi gọn sau gáy, vài sợi tóc mai lòa xòa trước trán. Trên tay bà ấy là một cây bút chì và cuốn sổ tay phác thảo ý tưởng cho buổi stream sắp tới.

Nhìn thấy tôi bước vào một mình với chiếc túi xách to, Kson chớp chớp mắt kinh ngạc, rồi khóe môi lập tức nhếch lên thành nụ cười ranh mãnh: ``Ủa, Kanatan? Không báo trước một tiếng gì hết vậy? Lại còn đi lẻ một mình, không dẫn theo bà Cừu hay công chúa kẹo ngọt sang à?''

``Đi một mình không được à?'' tôi bước tới, tháo giày đặt ngay ngắn bên cửa rồi ngồi bệt xuống thảm đối diện bà ấy. ``Hay là Tổng trưởng bận rộn quá, không rảnh tiếp khách quý thế này?''

``Nói linh tinh, khách quý như bà thì ngày nào tôi cũng tiếp!'' Kson cười khà khà, tiện tay đẩy đĩa bánh quy bơ sang trước mặt tôi. ``Sang đây có việc gì đột xuất thế? Trông mặt nghiêm túc chưa kìa.''

Tôi không trả lời ngay. Tôi nhẹ nhàng kéo khóa chiếc túi vải bố, cẩn thận nhấc cuốn album ảnh bọc nhung xanh dày cộp đặt lên chiếc bàn gỗ nhỏ.

Kson nhìn thấy cuốn album, nụ cười trên môi bà ấy bỗng khựng lại.

Đôi mắt tím của cô bạn mở to. Ánh nhìn của bà ấy dừng lại thật lâu trên góc bìa nhung sờn cũ, nơi có dòng chữ nắn nót do chính tay năm đứa chúng tôi ký tặng nhau vào mùa xuân năm 2020: \textit{holoForce - Mãi mãi bên nhau}.

Căn phòng bỗng chốc chìm vào một khoảng lặng êm đềm. Chỉ có tiếng gió hè khẽ lay động tấm rèm cửa màu trắng ngà.

``Cuốn này\dots'' Kson khẽ cất giọng, âm thanh trầm xuống một nhịp. ``Bà vẫn còn giữ à?''

``Tất nhiên là giữ chứ,'' tôi mỉm cười, vỗ nhẹ lên mặt bìa. ``Tôi đâu có ngốc đến mức vứt bỏ những thứ quý giá nhất đời mình. Nào, lại đây xem với tôi đi.''

Kson im lặng chừng vài giây, rồi bà ấy chậm rãi ngồi sát lại bên cạnh tôi. Hai bờ vai chạm nhẹ vào nhau, mang lại một cảm giác thân thuộc đến nao lòng, như thể khoảng thời gian một năm cách biệt kia chưa từng tồn tại.

Tôi lật mở trang đầu tiên.

Một tấm ảnh polaroid bốn góc dán băng dính hoa hiện ra: năm cô gái đứng chen chúc trong phòng tập vũ đạo ngầm, mặt mũi lấm lem mồ hôi, tóc tai bù xù nhưng nụ cười thì rạng rỡ như ánh mặt trời. Đó là tuần lễ tập dượt đầu tiên trước ngày ra mắt chính thức của Thế hệ thứ Tư.

``Trời đất ơi\dots'' Kson bật cười thành tiếng, đưa ngón tay chỉ vào góc ảnh. ``Nhìn mặt tôi hồi đó ngố chưa kìa! Tóc tai chẳng đâu vào đâu, lại còn mặc cái áo ba lỗ rách nách nữa chứ!''

``Bà thì lúc nào chẳng ngố,'' tôi trêu lại, nhưng ngón tay tôi cũng khẽ chạm vào góc áo của hình bóng Kson trong ảnh. ``Nhưng hồi đấy bà là người năng nổ nhất đấy. Ngày nào tập xong cũng lôi cả bọn đi ăn mì ramen, rồi ngồi thao thao bất tuyệt về ước mơ đưa cả nhóm vươn ra thế giới.''

Chúng tôi lật sang trang tiếp theo.

Tấm ảnh chụp trong căn hộ chung của tôi và bà ấy: một đêm mùa đông lạnh buốt, hai đứa ngồi cuộn tròn trong một chiếc chăn bông to đùng, cùng nhau ăn một nồi lẩu kim chi bốc khói nghi ngút. Bên cạnh là những bức ảnh chụp những buổi off-collab đầy hỗn loạn, những lần tập hát đến khản cả giọng, và cả khoảnh khắc năm đứa ôm nhau khóc nức nở sau sân khấu của buổi hòa nhạc chung đầu tiên.

Từng trang, từng trang ảnh lật qua, như một cuốn phim quay chậm đưa chúng tôi trở về những tháng năm thanh xuân rực rỡ nhất.

Khi lật tới trang cuối cùng - nơi chỉ còn lại một tấm ảnh chụp vội vào buổi sáng ngày 1 tháng 7 năm 2021, khoảnh khắc bà ấy thu dọn hành lý rời khỏi căn hộ chung - bàn tay tôi bỗng run rẩy.

Tôi dừng lại, ngước nhìn sườn mặt của Kson.

Dưới ánh nắng ban chiều rọi qua khung kính, bờ mi của bà ấy khẽ chớp, khóe mắt ánh lên những vệt sáng lấp lánh như ngấn nước.

``Bà biết không, Kson\dots'' giọng tôi nghẹn lại, từng chữ thốt ra từ tận đáy lòng sâu kín nhất. ``Suốt một năm qua, tôi đã giấu cuốn album này dưới đáy tủ. Tôi từng nghĩ mình sẽ không bao giờ có đủ can đảm để mở nó ra lần nữa. Tôi sợ rằng nếu nhìn vào những bức ảnh này, tôi sẽ không chịu nổi nỗi đau mất bà\dots\ Tôi từng nghĩ chúng ta sẽ vĩnh viễn không bao giờ được ngồi cạnh nhau lật từng trang ký ức này nữa.''

Kson không ngắt lời tôi. Bà ấy lặng lẽ vươn tay sang, bàn tay thô ráp, ấm áp siết nhẹ lấy mu bàn tay của tôi.

``Nhưng hôm nay ngồi ở đây, nhìn thấy bà mỉm cười thế này\dots\ tôi mới hiểu ra một điều,'' tôi hít một hơi thật sâu, mỉm cười qua làn nước mắt mỏng manh. ``Ký ức không mất đi đâu cả. Dù thế giới có đổi thay, chúng ta vẫn ở đây, vẫn là những người bạn thân nhất của nhau.''

Kson hít một hơi dài, ngửa đầu nhìn lên trần nhà để ngăn không cho giọt nước mắt rơi xuống. Khóe môi cô bạn cong lên một nụ cười rạng rỡ và dịu dàng lạ kỳ: ``Ừ. Ký ức này là của chúng ta, không ai có thể tước đoạt nó được cả, Kanatan ạ.''

Bà ấy khép cuốn album lại, nhẹ nhàng đặt tay lên ngực tôi: ``Cảm ơn bà vì đã giữ gìn nó cẩn thận như thế. Giờ thì cất nó đi thôi, hôm nay trời đẹp thế này, ngồi trong phòng mãi phí cả buổi chiều. Đi dạo với tôi một vòng không?''

``Đi chứ!'' tôi quẹt nhanh khóe mắt, mỉm cười gật đầu.

\ct

Phường Anwa về chiều khoác lên mình một vẻ đẹp yên bình hiếm thấy dù là một thành phố hiện đại.

Dù rằng nơi này cũng nườm nượp như ở Shibuya hay Shinjuku, nhưng lúc này, nơi đây chỉ có tiếng gió xào xạc qua những tán cây phong cổ thụ và tiếng nước róc rách của con sông chảy uốn lượn.

Tôi và Kson đi song song bên nhau trên con dốc lát đá mộc mạc. Cả hai đều diện đồ giản dị, đeo khẩu trang che nửa mặt để tận hưởng bầu không khí tự do ngoài trời.

``Dạo này công việc ở bên công ty thế nào rồi?'' Kson đút hai tay vào túi quần, cất tiếng hỏi bâng quơ. ``Nghe bảo bà sắp có dự án solo live lớn lắm đúng không?''

Tôi gật đầu, trong mắt ánh lên niềm say mê: ``Ừ, tôi đang chuẩn bị cho liveshow cá nhân đầu tiên. Có nhiều bài hát tôi tự viết lời, muốn truyền tải trọn vẹn thông điệp của một Amane Kanata độc lập, vững vàng trên đôi chân của chính mình. Lịch tập dày đặc lắm, nhưng tôi không thấy mệt chút nào cả.''

``Tốt lắm,'' Kson mỉm cười, giọng tràn đầy sự tự hào. ``Thiên sứ nhỏ ngày nào hay khóc nhè giờ đã trưởng thành thành một idol thực thụ rồi đấy nhỉ. Tôi tự hào về bà lắm, Kanatan.''

``Còn bà thì sao?'' tôi quay sang nhìn cô bạn. ``Kế hoạch của Tổng trưởng cho tương lai là gì?''

Kson dừng bước bên lan can đá của chiếc cầu nhỏ bắc qua con kênh, phóng tầm mắt nhìn về phía những ngọn đồi xanh biếc xa xa: ``Tôi à\dots\ Tôi muốn xây dựng một mái nhà cho riêng mình. Một cộng đồng VTuber độc lập, nơi mọi người có thể tự do sáng tạo, được là chính mình mà không phải chịu sự trói buộc của những quy chuẩn khắt khe hay những bản hợp đồng ngột ngạt. Tôi muốn chứng minh cho mọi người thấy rằng, ngay cả khi không có một tập đoàn lớn hậu thuẫn, người ta vẫn có thể sống tốt, cống hiến hết mình và mang lại tiếng cười cho khán giả bằng chính thực lực của mình.''

Nghe những lời đầy nhiệt huyết ấy, tôi ngắm nhìn bóng lưng vững chãi của Kson dưới ánh ráng chiều.

Hình ảnh của bà ấy lúc này thật kỳ diệu.

Trước đây, khi còn mang danh nghĩa một con Rồng thần thoại, bà ấy dang rộng đôi cánh rực rỡ để che chở cho chúng tôi trước giông bão. Giờ đây, khi cởi bỏ lớp vảy rồng và những hào quang cũ, bà ấy lại đang tự tay tạo dựng một đôi cánh mới—một đôi cánh tự do, kiên cường và đầy bản lĩnh.

Còn tôi, tôi vẫn đang tiếp tục dang rộng đôi cánh thiên sứ của mình trên bầu trời idol lấp lánh ánh đèn sân khấu.

Hai chúng tôi giờ đây đang bước đi trên hai ngả đường hoàn toàn khác biệt. Một người ở lại với thế giới thần tượng chính quy, một người tự do sải bước trên con đường độc lập. Thế nhưng, mục đích cuối cùng của cả hai lại chưa từng rẽ hướng: chúng tôi đều muốn mang lại niềm hy vọng, sự ấm áp và tiếng cười cho những người luôn dõi theo mình.

Hai đôi cánh khác nhau, nhưng cùng chung một lý tưởng sống.

\ct

Chúng tôi dừng chân tại một công viên nhỏ trên lòng sông, nơi có thể ngắm trọn vẹn toàn cảnh thị trấn Anwa đang dần chìm trong sắc tím hoàng hôn rực rỡ.

Gió thổi về lồng lộng, làm tung bay mái tóc ngắn của tôi và tà áo của Kson.

Tôi đứng tựa vào lan can gỗ, nhìn sườn mặt góc cạnh của người bạn thân. Cảm xúc nghẹn ngào từ cuốn album ban nãy lại một lần nữa trào dâng. Bất giác, tôi bước lại gần một bước, đưa tay ra nắm chặt lấy bàn tay thô ráp của bà ấy.

Kson thoáng giật mình, quay sang nhìn tôi với ánh mắt ngơ ngác: ``Gì thế, Kanatan?''

Tôi siết chặt bàn tay bà ấy bằng tất cả sức mạnh năm mươi cân của mình—nhưng không phải để đo lực tay, mà là để gửi gắm trọn vẹn sự kiên định của trái tim:

``Kson này. Dù tương lai có ra sao, dù bà có đi xa đến đâu và làm bất cứ điều gì\dots\ tôi cũng sẽ không bao giờ để mất bà thêm một lần nào nữa.''

Tôi nhìn thẳng vào đôi mắt tím thẳm của cô bạn, từng chữ rõ ràng và dứt khoát: ``Đây là lời hứa của tôi với bà. Cho dù hai đứa đi trên hai con đường khác nhau, chỉ cần bà cần, tôi sẽ luôn có mặt. Mãi mãi là như vậy.''

Kson đứng sững sờ. Đôi mắt tím của bà ấy dao động mãnh liệt trước ánh nhìn kiên định của tôi.

Rồi, một nụ cười ấm áp, rạng ngời và dịu dàng nhất từ trước đến nay từ từ nở trên khóe môi cô bạn. Bà ấy lật bàn tay lại, siết chặt lấy những ngón tay nhỏ nhắn của tôi:

``Ừ. Tôi biết rồi, đồ ngốc ạ. Tôi cũng hứa với bà như thế.''

Hai bàn tay đan chặt vào nhau giữa hoàng hôn lộng gió, như một lời thề son sắt giữa hai linh hồn tri kỷ đã cùng nhau đi qua giông bão cuộc đời.

\ct

Bảy giờ ba mươi phút tối.

Ánh đèn vàng của ga Kyuukami đã bật sáng lung linh. Tiếng chuông báo tàu chuẩn bị vào ga vang lên từng hồi giòn giã.

Kson tiễn tôi ra tận cửa soát vé. Trên vai bà ấy vẫn là chiếc áo khoác gió mỏng, hai tay đút túi quần, điệu bộ thong dong quen thuộc.

``Về đến nhà nhớ nhắn tin đấy nhé,'' Kson dặn dò. ``Đừng có mải tập hát quá mà bỏ bữa tối, biết chưa?''

``Biết rồi, bà lo cho tôi như mẹ già ấy!'' tôi bật cười, giơ chiếc vé tàu trên tay lên vẫy vẫy.

Tôi quay lưng chuẩn bị bước qua cổng từ, thì bất chợt giọng nói của Kson từ phía sau vang lên, trầm ấm và chân thành hơn bao giờ hết:

``Này Kanata.''

Tôi dừng bước, ngoái đầu lại.

Dưới ánh đèn vàng ấm áp của nhà ga, Kson đứng đó, mái tóc tím than khẽ đung đưa theo làn gió đêm. Cô bạn mỉm cười, đôi mắt sáng lấp lánh:

``Cảm ơn bà\dots\ vì đã đến tìm tôi.''

Cổ họng tôi nghẹn lại, một cảm giác ngọt ngào lan tỏa khắp lồng ngực. Tôi nở nụ cười rạng rỡ nhất của một thiên sứ, cất giọng thật trong trẻo:

``Tôi đã muốn nói câu đó từ lâu rồi. Gặp lại sau nhé, Tổng trưởng!''

Đoàn tàu từ từ lăn bánh, mang theo tôi trở về Tokyo dưới bầu trời đêm đầy sao. Qua ô cửa kính, tôi nhìn thấy bóng dáng Kson đứng vẫy tay trên sân ga xa dần, xa dần.

Lòng tôi lúc này nhẹ nhõm và bình yên đến lạ kỳ.

Hai đôi cánh của chúng tôi có thể đang bay về hai phương trời khác biệt, nhưng sợi dây liên kết vô hình giữa hai trái tim sẽ mãi mãi là ngọn hải đăng soi sáng con đường phía trước.

%==========================================TRUYỆN PHỤ=========================================
\omk

\pov{-}

Khoảng tám giờ tối hôm đó, tại phòng khách tầng năm nhà Hikawa.

Kson lững thững bước vào nhà, cởi đôi dép lê để ngay ngắn trên kệ rồi đi thẳng vào phòng bếp định tìm nước uống. Khóe mắt cô nàng vẫn còn hơi ửng đỏ, dù nụ cười trên môi thì không giấu nổi vẻ thỏa mãn.

Vừa bước qua cửa bếp, một giọng nói lanh lảnh đầy vẻ tinh quái đã bất ngờ vang lên từ chiếc bàn ăn: ``Úi chà chà! Nhìn ai kìa!''

Ryuuhou đang ngồi gác chân lên ghế, một tay cầm que kem dưa hấu cắn dở, hai mắt sáng rực như đèn pha săm soi nhìn chằm chằm vào mặt Kson: ``Nè Kson-san! Chị vừa khóc nhè đấy à? Hai mắt đỏ hoe như mắt thỏ thế kia kìa!''

Kson giật bắn mình, vội vàng đưa mu bàn tay lên quẹt mắt, hắng giọng gắt gỏng: ``Khóc lóc cái gì chứ nhóc con! Gió to ngoài đường bụi bay vào mắt thôi!''

Ryuuhou cười khẩy một tiếng đầy mỉa mai, chỉ que kem ra ngoài cửa sổ: ``Chị lừa con nít đấy à? Tối nay gió lặng như tờ, với lại ở cả cái phường Anwa này cây xanh rợp bóng, vườn tược khắp nơi đào đâu ra bụi cho chị bay vào mắt? Thừa nhận là chia tay Kanata-san xúc động quá khóc nhè đi cho rồi!''

Kson khựng lại một nhịp, đôi má lập tức đỏ bừng lên vì ngượng. Cô nàng bước nhanh tới, vung bàn tay dài ngoẵng cốc nhẹ một cái *CỐC* lên đỉnh đầu cô bé tinh nghịch: ``\dots Im mồm đi, nhóc con! Ăn kem cho nghẹn họng bây giờ!''

``Ai da!'' Ryuuhou ôm đầu làu bàu nhưng miệng vẫn cười toe toét: ``Bị nói trúng tim đen nên thẹn quá hóa giận chứ gì! Để em đi mách Onii với Kaiho-nee là Tổng trưởng giang hồ khóc nhè nhé!''

``Mày dám mách thử xem, mai tao bắt mày sang xưởng cơ khí phụ vác đồ mười chuyến cho chừa!''

Tiếng cãi cọ chí choé rộn rã lại vang lên khắp gian bếp tầng năm, làm tan biến đi mọi sự lắng đọng vừa rồi, mang lại một buổi tối tràn ngập tiếng cười bình dị dưới mái nhà Hikawa.

\end{document}
